\begin{annexitem}{annex:p2.tr.1}{thm:p02-translation-composition}{The finite obstruction}
\begin{proof}
Let \(\mathbf V\) be the distinguished universe symbol. Put
\[
 L=\{\mathbf V,P,c\},\quad
 K=\{\mathbf V,U,P,c\},\quad
 J=\{\mathbf V,D,W,P,c\},
\]
where all displayed predicates are unary and \(c\) is a constant.
Define zero-term interpretations
\[
 \begin{array}{c|ccc}
 \alpha:L\to K&\mathbf V\mapsto U&P\mapsto P&c\mapsto c\\
 \beta:K\to J&\mathbf V\mapsto D&U\mapsto W&P\mapsto P,\ c\mapsto c.
 \end{array}
\]
Both obey the source's symbol and arity requirements. Componentwise
composition sends \(\mathbf V\) to \(W\), and fixes \(P,c\).

Take the atomic sentence \(\varphi=P(c)\). Its least language is \(L\).
Indeed truth depends on both \(P\) and \(c\), and every language contains
\(\mathbf V\). To justify this occurrence argument directly, inclusions
are simple interpretations and hence always invertible. Functoriality
with the inclusion of a least language, together with the Translation
Axiom, implies that two structures agreeing on that least language
agree on the sentence's truth. Varying \(P\) or \(c\) on a fixed
two-element universe changes the truth of \(P(c)\), so neither can be
omitted from its least language.

For any \(K\)-structure \(B\), the guarded Translation Axiom gives
\[
 B\models\varphi^\alpha
 \quad\Longleftrightarrow\quad c^B\in U^B\ \hbox{ and }\ c^B\in P^B.
\]
The least language of \(\varphi^\alpha\) is therefore exactly \(K\).
For completeness, on the universe \(\{0,1\}\), holding
\(U=P=\{1\}\) fixed and changing \(c\) from 1 to 0 changes truth;
holding \(U=\{1\},c=1\) fixed and changing \(P\) from \(\{1\}\) to
\(\varnothing\) changes truth; and holding \(P=\{1\},c=1\) fixed
and changing \(U\) from \(\{1\}\) to \(\{0\}\) changes truth.
The occurrence argument excludes omission of each of \(c,P,U\).
Thus the Translation Axiom applies to \(\beta\) with precisely its
displayed source language, not an assumed larger support.

Now let \(N\) be the \(J\)-structure with
\[
 |N|=\{0,1\},\qquad D^N=\{0\},\qquad
 W^N=P^N=\{1\},\qquad c^N=1.
\]
It is not \(\beta\)-invertible: its selected domain is \(\{0\}\), which
does not contain the value of \(c\). Therefore
\(N\not\models(\varphi^\alpha)^\beta\).
In contrast it is \(\beta\alpha\)-invertible. That pull-back has
universe \(\{1\}\), predicate \(\{1\}\), and constant 1, so it
satisfies \(P(c)\). Consequently
\(N\models\varphi^{\beta\alpha}\). Functoriality would identify these
two sentences, a contradiction. Every selected nonempty domain in the
example is explicit; the contradiction does not rely on an empty-domain
convention or on a quantifier axiom.
\end{proof}
\end{annexitem}

\begin{annexitem}{annex:p2.tr.2}{thm:p02-guarded-fo-translation}{The guarded operation}
\begin{proof}
A target structure satisfies the finite guard exactly when every
interpreted source function maps tuples of the selected domain back to
that domain, and each interpreted source constant belongs to it. These
are precisely the pull-back's totality requirements. Relation symbols
are restricted to the domain and introduce no additional totality
condition. The statement remains correct for an empty selected domain:
positive-arity closure conditions are vacuous, while a constant's
condition fails.

Assume the guard. Induction on source terms shows that their values in
the pull-back are their translated values in the target: constants are
the zero-arity case, and a function application uses its arity-raised
symbol at the fixed parameter tuple. Closure ensures that all those
values remain in the selected domain. Atomic formulas and equality
therefore have the same truth values. Boolean operations preserve this
comparison. An existential quantifier in the source ranges exactly
over elements satisfying \(D_\alpha\) in the target, so its translation
is \(\exists x\,(D_\alpha(x)\land\widehat\chi(x))\). Universals translate
as \(\forall x\,(D_\alpha(x)\to\widehat\chi(x))\). Structural induction
proves truth preservation. Conjoining the guard makes the translation
false when the pull-back is not a structure, giving the exact guarded
condition.

For a simple zero-term interpretation the universe is unchanged and
all target functions are already total. No guard is needed. Ordinary
symbol substitution preserves identities and composition, by induction
on terms and formulas, including capture-free bound variables. Thus
these symbol changes do act functorially.

Finally apply the individually proved guarded operation successively
for a sequence of interpretations. At each step it describes the
pull-back required at that step, with the occurring intermediate symbols
and their guard. Repeated application of the established equivalence
gives the sequential truth condition. The earlier counterexample shows
why discarding an intermediate guard is not sound. Even when no functions
occur, an intermediate domain can matter: if a first relativization is
to \(U\) and a second one selects \(D\), with \(U\) renamed to \(W\),
the sequential domain is \(D\cap W\), whereas the componentwise
symbol map selects only \(W\). Retaining the successive operations
preserves their intersection. No unsupported full-functor assertion
has been used.
\end{proof}
\end{annexitem}
